\documentclass[11pt]{article}

% DOCUMENT SETUP
\usepackage[utf8]{inputenc}
\usepackage[margin=1in]{geometry}
\usepackage{fullpage}
\setlength{\parindent}{0pt}

% FONTS & TEXT
\usepackage{soul}
\usepackage{slantsc}
\usepackage{bold-extra}
\usepackage{fontawesome5}

% FANCY HEADERS (optional - uncomment to use)
%\usepackage{fancyhdr}
%\pagestyle{fancy}
%\rhead{Event/Conference}
%\lhead{Date}
%\fancyhfoffset{0in}
% MATH & SYMBOLS
\usepackage{amsmath}
\usepackage{amssymb}
\usepackage{amsfonts}

% GRAPHICS & FIGURES
\usepackage[demo]{graphicx}
\usepackage{subcaption}
\usepackage{float}

% TABLES
\usepackage{multicol}
\usepackage{multirow}
\usepackage{array}
\usepackage{booktabs}

% COLORS
\usepackage{color}
\usepackage[dvipsnames]{xcolor}
\definecolor{pastelgreen}{RGB}{193, 225, 193}
\definecolor{peach}{RGB}{255, 229, 180}
\definecolor{airforceblue}{rgb}{0.36, 0.54, 0.66}
\definecolor{camouflagegreen}{rgb}{0.47, 0.53, 0.42}
\definecolor{darkcyan}{rgb}{0.0, 0.55, 0.55}

% TIKZ & PLOTTING
\usepackage{tikz}
\usepackage{tikz-qtree}
\usepackage{tikz-qtree-compat}
\usetikzlibrary{arrows, positioning, shadows}
\usetikzlibrary{arrows.meta}
\usetikzlibrary{fit}
\tikzset{every tree node/.style={align=center,anchor=north}}

\usepackage{pgf-pie}
\usepackage{pgfplots}
\usepackage{pgfplotstable}
\pgfplotsset{compat=1.17}

% BOXES & FRAMES
\usepackage{tcolorbox}

% LINGUISTICS
\usepackage{tipa}
\usepackage{gb4e}
\usepackage{linguex}

% TABLE OF CONTENTS & SECTIONS
\usepackage{titletoc, tocloft}
\usepackage{etoolbox}
\usepackage{titlesec}

% Section formatting
\renewcommand\thesection{\Roman{section}}
\titleformat{\section}
  {\large\bf\scshape}{\thesection}{0.2in}{}
\renewcommand\cftsecfont{\bf\MakeUppercase}
\setlength{\cftsecnumwidth}{3.5em}

\makeatletter
\patchcmd{\l@section}
    {\cftsecfont #1}
    {\cftsecfont {#1}}
    {}
    {}
\makeatother

% BIBLIOGRAPHY
\usepackage[square,numbers]{natbib}

% FOOTNOTES
\usepackage[symbol]{footmisc}

% HYPERLINKS (load last)
\usepackage{hyperref}
\hypersetup{
    colorlinks=true,
    linkcolor=purple,
    citecolor=camouflagegreen,
    urlcolor=darkcyan,
}

% CUSTOM COMMANDS
\newcommand{\myframe}[2]{
    \begin{tcolorbox}[colback=white,colframe=black,sharp corners, title={#1}]
        #2
    \end{tcolorbox}
}

% DOCUMENT START
\begin{document}
%TC:ignore 
\begin{center}
   \Large{\textsc{\textbf{[Sycophancy as a Function of Context Topic]}}}\\
   \vspace{0.1in}
   \normalsize
   March 2026 \\
\end{center}
%TC:endignore 
\hrule
\vspace{0.2cm}

\begin{multicols}{2}
   \textbf{\textsc{Introduction.}} Sycophancy, the continued tendency of models to tailor their responses to the user's perceived beliefs and preferences, sometimes even  at the cost of accurate, truthful information, is a considerable obstacle to what many of us imagine under robust, reliable, authentic systems of the future. While this at best obsequious and possibly even ingenuine flattery underlying models' reactions to the surrounding world is often viewed as a capability defect, it has become increasingly more evident that sycophancy is, first and foremost, a major safety issue. That said, a notable complication in tackling sycophancy is insufficient goal specificity: sycophancy may not be an undesirable trait in all contexts and it may be problematic in different ways across different contexts. This study aims to bring more clarity into how sycophancy manifests within the conversation context and inform thus more targeted, efficient interventions.\\

   \textbf{\textsc{Lit Review.}}

\end{multicols}



\newpage
\textbf{The Sycophancy Problem.}  

Sycophancy has been mostly viewed as a capability defect: sycophancy comes hand-in-hand with something of a truth-indifference, preventing AI from becoming a [decisive] integrated part of society. However, evidence has been increasingly suggesting that sycophancy is also a major, and often overlooked, safety issue. [...] This effect is further reinforced by the fact that models seem to put a unique weight on the user's mental model of the world, as suggested by the fact that prompts including an ``I believe'' statement trigger sycophancy a lot more than prompts including ``They believe'' statements \citep{wang2025}. [In other words, what sycophancy creates at scale may not be an existential risk, but what it creates while it scales may be.]

[What the future will look like \& why this research is necessary for it]

Sycophancy also happens in STEM topics: \citep{petrov2025}

Another question is how to approach sycophancy elimination. \citep{noshin2026} argues that some contexts require sycophantic responses. At the same time, whether or not a model is sycophantic currently seems to depend on another axis: subjectivity. [...] Sycophancy occurs significantly more on more subjective tasks thank objective tasks \citep{ranaldi2025}. 

Sycophancy is a non-trivial problem because it's somewhat difficult to measure.

Interaction context triggers sycophancy, but in what ways remains poorly understood \citep{jain2025}

[Where sycophancy comes from.]
Humans tend to prefer convincingly-written and sycophantic responses over correct ones, resulting in overly flattering, not necessarily accurate RLHF models \citep{sharma2025}

In the future, AI could discriminate against humans \citep{laurito2025}

Relatively small, constant poisoned data points in any training sample cause predictable output \citep{souly2025}, which could relate to sycophancy in that we can observe whether this finding transfers to the quantity of data points by topic impacts sycophancy in output.

Consequences of AI include dependency \citep{cheng2026}

\newpage
\bibliographystyle{plain}
\bibliography{lib}
\nocite{*}

\end{document}